\section{Interpretation and Limits}

The formal state has five coupled structural observables and a temporal
transition layer. Topology concerns qualitative organization such as adjacency
and connectedness. Geometry concerns metric information. Relations concern
typed entities, generator morphisms, and composition. Filtration and mass add
ordered and measure-theoretic observables on the same carrier. Dynamics concerns
action-conditioned tracked transitions, histories, turnover, interventions, and
rollouts. The common tagged carrier and the compatibility axiom couple these
layers without identifying them.

The formal results establish conditional statements about explicitly specified
objects. They do not prove that these observables are necessary or sufficient
for intelligence, that human cognition literally uses this state object, or
that an unconstrained neural network will discover it. They also do not prove
that a structurally specified model is universally superior. Those statements
require separate cognitive and empirical validation.

The principal design boundary is therefore deliberate: a structural claim must
name its carrier, equivalence, observables, transformations, and discrepancy
before it can be evaluated. A scalar predictive loss may be useful, but it
cannot by itself determine structural equivalence or transition preservation.
Conversely, a structural discrepancy is meaningful only after its task-relevant
alignment and invariants are specified.

The formalism takes the entity set, its typing, and its cross-time identity as
given data of the state. It does not supply a procedure for discovering
entities from raw observation or for maintaining identity under perception
error. Those stages precede everything specified here.

The companion empirical papers test bounded consequences of this framework on
finite mechanism families. Those experiments should be read as validation of
the stated consequences, not as validation of the entire theory or of
real-world cognition.
